\chapter{Методика построения, исследования и аппроксимации многообразия решений энергетически оптимальных задач встречи}

% =========================================================
% Рабочий каркас главы 1.
% Это НЕ финальная редакция главы, а удобная структура для
% аккуратного переноса материалов из ВКР и презентации.
%
% Принцип заполнения:
% - сначала переносится опорный текст из ВКР;
% - затем он дополняется и перестраивается по логике презентации;
% - служебные комментарии и указания на источники удаляются только
%   на позднем этапе шлифовки.
% =========================================================

% Общая опора:
% ВКР: главы 1--2, с. 12--22
% Презентация: слайды 11--18

\section{Энергетически оптимальная задача встречи}
% Опора:
% ВКР: §1.1, с. 12--13
% Презентация: слайды 13--14
%
% Здесь последовательно ввести:
% - уравнения движения КА в модели двух тел;
% - реактивное ускорение и допущение об идеально регулируемой тяге;
% - квадратичный функционал;
% - связь функционала с расходом рабочего тела;
% - граничные условия задачи встречи.
%
% Удобный порядок переноса:
% 1) базовые уравнения движения;
% 2) ограничения ИР-модели;
% 3) функционал;
% 4) интерпретация через расход топлива;
% 5) задача встречи.
%
% Черновой текст:

\section{Необходимые условия оптимальности и параметры оптимального управления}
% Опора:
% ВКР: §1.2, с. 14
% Презентация: слайд 15
%
% Здесь собрать:
% - функцию Гамильтона--Понтрягина;
% - условие максимума;
% - выражение оптимального управления;
% - уравнения для сопряжённых переменных;
% - объяснение, почему начальные значения сопряжённых переменных
%   задают параметры оптимального управления вдоль траектории.
%
% Важно:
% - не просто повторить формулы из ВКР,
% - а явно связать их с дальнейшей задачей аппроксимации параметров
%   начального приближения.
%
% Черновой текст:

\section{Переход к безразмерным переменным}
% Опора:
% ВКР: §2.1, с. 15
% Презентация: слайды 13--14
%
% Здесь ввести:
% - нормировку радиуса начальной орбиты;
% - нормировку гравитационного параметра;
% - единицы DU, VU, TU, AU;
% - масштабирование сопряжённых переменных.
%
% Надо оставить акцент:
% - безразмерная форма нужна не только для удобства записи,
% - но и для унификации семейства рассматриваемых перелётов.
%
% Черновой текст:

\section{Выбор граничных условий и построение модельного многообразия}
% Опора:
% ВКР: §2.2, с. 16--17
% Презентация: слайды 7, 14
%
% Здесь объяснить:
% - почему реальные массивы траекторий неудобны для прямой аппроксимации;
% - почему вводятся виртуальные планеты;
% - как задаются синодическая фаза, угловая / витковая дальность,
%   отношение радиусов;
% - почему именно такое пространство параметров становится
%   центральным объектом исследования.
%
% Желательно явно использовать термин:
% "модельное многообразие".
%
% Черновой текст:

\section{Метод построения многообразия на основе дискретного продолжения}
% Опора:
% ВКР: §2.3, с. 18--20
% Презентация: слайд 16
%
% Здесь последовательно описать:
% - получение базового решения;
% - распространение решения по координатной сетке;
% - вычислительные барьеры;
% - исключаемые области;
% - смысл дискретного продолжения как практического способа
%   построения гладкого набора решений.
%
% Важно:
% - в этой версии главы не надо перегружать раздел техническими
%   подробностями MATLAB-реализации;
% - акцент должен быть на методической идее.
%
% Черновой текст:

\section{Методика исследования и аппроксимации многообразия}
% Опора:
% ВКР: §2.3--2.4, с. 18--22
% Презентация: слайды 12, 17--18
%
% Здесь удобно дать именно общий цикл:
% 1) построение многообразия;
% 2) выбор параметризации;
% 3) аппроксимация;
% 4) исследование аппроксимаций;
% 5) репараметризация при необходимости;
% 6) уточнение настроек символьной регрессии;
% 7) получение конечных формул.
%
% Этот раздел важен как "скелет" всей диссертации:
% именно на него потом будут опираться главы 2--4.
%
% Черновой текст:

\section{Символьная регрессия как инструмент получения конечных формул}
% Опора:
% ВКР: §2.4, с. 21--22
% Презентация: слайды 6, 17
%
% Здесь нужно:
% - кратко мотивировать выбор символьной регрессии;
% - показать её отличие от полиномиальной / спектральной аппроксимации
%   и от менее интерпретируемых моделей;
% - зафиксировать, что метод используется для автоматического поиска
%   структуры конечных формул и оценки коэффициентов;
% - при необходимости кратко упомянуть предварительное сравнение
%   альтернативных подходов.
%
% Черновой текст:

\section{Критерии отбора аппроксимационных формул и оценка качества аппроксимаций}
% Опора:
% ВКР: §2.4, с. 21--22
% Презентация: слайды 17, 21, 24
%
% Здесь собрать критерии, которые дальше повторяются по всей работе:
% - отсутствие сингулярностей в области определения;
% - ошибка на отложенной выборке;
% - корректное поведение в предельных случаях;
% - физическая интерпретируемость;
% - пригодность для практического применения в алгоритме начального
%   приближения.
%
% Здесь же удобно заранее зафиксировать мысль, что в работе отбираются
% не просто "наиболее точные" формулы, а формулы, сочетающие точность,
% устойчивость и интерпретируемость.
%
% Черновой текст:

\section*{Выводы к главе 1}
% Опора:
% ВКР: главы 1--2, с. 12--22
% Презентация: слайд 18
%
% Здесь потом кратко зафиксировать:
% 1. постановку энергетически оптимальной задачи встречи;
% 2. параметры оптимального управления через начальные сопряжённые
%    переменные;
% 3. способ построения модельного многообразия;
% 4. общую методику исследования и аппроксимации;
% 5. критерии получения конечных формул.
%
% Черновой текст:
